\chapter{Investigação Académica e Relatórios Técnicos}
\label{ch:investigacao}

\section{O que a IA faz bem (e mal) em contexto académico}

Usar LLM em investigação e ensino superior exige escolhas explícitas sobre método e transparência. Pontos fortes:

\begin{itemize}[leftmargin=*]
\item ajuda a clarificar perguntas de investigação;
\item sugere estruturas para artigos, dissertações, relatórios;
\item melhora legibilidade de texto técnico;
\item gera checklists e grelhas de correção.
\end{itemize}

Pontos fracos (e perigosos):

\begin{itemize}[leftmargin=*]
\item inventa referências, DOIs e citações plausíveis mas falsas;
\item pode distorcer resultados de artigos se resumir mal;
\item facilita plágio e “trabalhos ocos” se usado para escrever tudo.
\end{itemize}

\begin{warningboxenv}{Se o modelo “inventa” artigos, a responsabilidade é tua}
Os modelos não têm uma lista fidedigna e estática de papers aprovados. Se pedes “cita 10 artigos sobre X”, ele pode inventar títulos e autores.
Nunca uses referências sugeridas por IA sem as verificares em bases reais (Google Scholar, bases de dados da biblioteca, etc.).
\end{warningboxenv}

\section{Refinar perguntas de investigação e objetivos}

Uma boa pergunta de investigação é:

\begin{itemize}[leftmargin=*]
\item clara;
\item focada (não tenta salvar o mundo inteiro);
\item operacionalizável (dá para recolher dados e responder de forma objetiva ou argumentada).
\end{itemize}

Padrão de uso da IA:

\begin{lstlisting}
Vou descrever o tema geral em que quero trabalhar
e uma primeira versão da minha pergunta de investigação.

Tema geral:
[texto]

Primeira tentativa de pergunta:
[texto]

Quero que:

critiques esta pergunta quanto a clareza,
foco e exequibilidade;

proponhas 3 versões alternativas mais focadas,
justificando o que muda em cada uma;

identifiques, para a versão que consideras
mais promissora, que tipo de dados ou fontes
seriam necessárias para responder.
\end{lstlisting}

\begin{promptlabbox}{Usar IA como ``orientador preliminar''}
\textbf{Objetivo e contexto:} Delimitar uma investigação fictícia ou pública, com tempo e dados disponíveis explícitos.

Escolhe um tema em que estejas a pensar fazer trabalho (e.g., segurança em redes, justiça, mobilidade urbana). Pede ao modelo para:

\begin{enumerate}[leftmargin=*]
\item ajudar a transformar esse tema em 3–5 perguntas de investigação distintas;
\item para cada uma, listar:
\begin{itemize}[leftmargin=*]
\item tipo de estudo provável (experimental, observacional, teórico),
\item fontes ou dados necessários,
\item principais dificuldades previsíveis.
\end{itemize}
\end{enumerate}

Depois discute as opções com um orientador humano. A IA ajuda a preparar o terreno, não a decidir por ti.

\textbf{Resultado esperado e validação:} Verifica se as perguntas são exequíveis e se dados e método respondem ao objetivo. Discute a escolha com um orientador, sem tratar sugestões de fontes como referências verificadas.

\textbf{Falhas possíveis:} Inventar acesso a amostras, formular causalidade sem desenho adequado ou aceitar uma metodologia só pelo vocabulário.
\end{promptlabbox}

\section{Estruturar artigos, dissertações e relatórios técnicos}

Modelos são bons a reproduzir estruturas padrão (IMRaD, relatórios técnicos com secções típicas, etc.). Em vez de pedires “escreve-me o relatório”, pede:

\begin{lstlisting}
Quero um esqueleto detalhado para um relatório
técnico sobre [tema], com cerca de [n] páginas.

Formato:

Índice proposto com capítulos e secções.

Para cada secção: 2-3 frases a indicar o que
deve ser tratado ali.

Indica também que tipo de figuras ou tabelas
fariam sentido em cada parte.
\end{lstlisting}

Depois, podes pedir ajuda secção a secção. Padrões úteis:

\subsection*{Introdução}

\begin{lstlisting}
Preciso de uma primeira versão de Introdução
(1-2 páginas) para um relatório sobre [tema].

Requisitos:

Contextualizar o problema sem encher chouriços.

Apresentar claramente objetivo(s) do trabalho.

Delimitar o escopo (o que fica de fora).

Não inventar dados nem referências específicas.

Depois desta versão, quero que proponhas
3-5 perguntas que devo responder para
refinar e personalizar o texto.
\end{lstlisting}

\subsection*{Metodologia}

\begin{lstlisting}
Abaixo descrevo a metodologia que realmente usei.

[texto teu, em bruto]

Quero que organizes isto numa secção de Metodologia
clara, com subsecções (ex: amostra, instrumentos,
procedimentos, análise de dados), sem inventar
nada que eu não tenha descrito.

Se achares que faltam detalhes importantes,
aponta-os no fim.
\end{lstlisting}

\section{Trabalhar com fontes: leitura, síntese e comparação}

Aqui a IA brilha se lhe deres material explícito. Em vez de pedires “resume a literatura sobre X”, faz:

\begin{itemize}[leftmargin=*]
\item recolhe tu os artigos relevantes (via bases académicas);
\item para cada artigo, copia o resumo e, se puderes, partes-chave;
\item pede ao modelo para:
\begin{itemize}[leftmargin=*]
\item sintetizar,
\item comparar,
\item extrair elementos estruturados.
\end{itemize}
\end{itemize}

Padrões:

\subsection*{Resumo estruturado de artigo}

\begin{lstlisting}
A partir do texto abaixo (abstract + excertos),
cria um resumo estruturado com:

Pergunta de investigação

Metodologia

Principais resultados

Limitações mencionadas

Ideias para trabalho futuro

[texto]
\end{lstlisting}

\subsection*{Comparar vários artigos}

\begin{lstlisting}
Vou colar resumos de 4 artigos sobre [tema].

Para cada um:

identifica pergunta, método, resultados em 2-3 frases.

Depois, faz uma síntese comparativa:

em que pontos concordam?

em que pontos divergem?

que lacunas parecem comuns?

[Abstract 1]
[Abstract 2]
[...]
\end{lstlisting}

\begin{checklistbox}{Uso disciplinado da IA na revisão de literatura}
Para cada pacote de artigos:

\begin{itemize}[leftmargin=*]
\item obtiveste os textos em fontes académicas legítimas?
\item usaste a IA apenas para ajudar a sintetizar e comparar?
\item mantiveste notas tuas (não apenas do modelo) sobre o que achas importante?
\item verificaste manualmente interpretações críticas?
\end{itemize}
\end{checklistbox}

\section{Escrita clara de texto técnico}

Textos académicos não têm de ser opacos. LLMs podem ajudar a:

\begin{itemize}[leftmargin=*]
\item simplificar frases sem perder rigor;
\item eliminar repetições;
\item adaptar o nível de formalidade.
\end{itemize}

Padrões típicos:

\subsection*{Clareza sem perder precisão}

\begin{lstlisting}
Reescreve o texto abaixo em português de Portugal,
mantendo o conteúdo técnico, mas:

frases mais curtas,

menos voz passiva,

evitar jargão desnecessário.

Se achares frases muito ambíguas, aponta-as e
propõe duas alternativas.

[texto]
\end{lstlisting}

\subsection*{Versão para diferentes públicos}

\begin{lstlisting}
A partir do texto técnico abaixo, produz:

Uma versão para estudantes de licenciatura
no mesmo curso (assume alguma base).

Uma versão para público geral interessado,
sem formação técnica.

Assinala claramente qual é qual.

[texto]
\end{lstlisting}

\begin{warningboxenv}{Não deixes a IA inventar resultados}
Pede ajuda na \emph{forma}, não no \emph{conteúdo} dos resultados.
Se pedes “escreve a secção de resultados” sem lhe dar os dados, o modelo pode fabricar uma continuação plausível; deve recusar preencher resultados sem dados.
Resultados e dados \emph{vêm das tuas análises}, não da imaginação da IA.
\end{warningboxenv}

\section{Ética, plágio e transparência}

Ponto sensível: é fácil escorregar para:

\begin{itemize}[leftmargin=*]
\item deixar a IA escrever metade do trabalho e fingir que é tudo teu;
\item copiar literalmente redações sugeridas;
\item usar referências inventadas;
\item violar regras explícitas da instituição.
\end{itemize}

Boas práticas:

\begin{itemize}[leftmargin=*]
\item Lê e respeita as políticas da tua universidade sobre uso de IA.
\item Usa a IA para:
\begin{itemize}[leftmargin=*]
\item brainstorming,
\item estruturação,
\item melhoria de redação,
\item síntese de materiais que tu próprio forneceste.
\end{itemize}
\item Não uses para:
\begin{itemize}[leftmargin=*]
\item escrever trabalhos inteiros “de raiz”,
\item inventar dados,
\item inventar citações.
\end{itemize}
\item Quando fizer sentido, declara o uso de IA (por exemplo, “foi usada ferramenta de IA generativa para auxílio na redação e revisão de texto, sob supervisão do autor”).
\end{itemize}

\begin{keyideabox}{A IA pode ajudar a pensar melhor, não a pensar menos}
Se usas o modelo para te poupar do esforço cognitivo essencial (ler fontes, entender métodos, estruturar argumentos), estás a sabotar o teu próprio crescimento académico.
Se o usas para te libertar de fricção superficial (formatação, algumas reescritas), podes focar mais no que interessa.
\end{keyideabox}

\section{Preparar-se para revisores, orientadores e bancas}

Finalmente, podes usar LLMs para simular o “lado de lá”:

\begin{itemize}[leftmargin=*]
\item revisores de revista,
\item orientadores exigentes,
\item membros de banca.
\end{itemize}

Padrões:

\subsection*{Simulação de revisor anónimo}

\begin{lstlisting}
Assume o papel de revisor anónimo de uma revista
científica na área de [área].

A partir do texto abaixo (Introdução + Metodologia):

Resume em 5-7 frases o que entendes ser o
contributo do trabalho.

Aponta 5 pontos fortes.

Aponta 5 pontos fracos ou dúvidas.

Sugere 3 melhorias específicas que aumentariam
a qualidade do artigo.

[texto]
\end{lstlisting}

\subsection*{Simular perguntas de banca}

\begin{lstlisting}
Assume que és membro de uma banca de dissertação
de mestrado em [área].

Tendo em conta o resumo do trabalho abaixo, gera
10 perguntas que poderias fazer ao candidato,
focadas em:

justificações de escolha metodológica,

limitações do estudo,

implicações práticas,

possíveis extensões futuras.

[resumo da dissertação]
\end{lstlisting}

\begin{promptlabbox}{Vacina anti-surpresa de banca}
\textbf{Objetivo e contexto:} Preparar defesa oral com base apenas no resumo e conclusões fornecidos.

Pega no resumo e conclusões do teu trabalho (ou de um trabalho fictício) e pede:

\begin{lstlisting}
Gera 15 perguntas difíceis e incómodas que
um orientador cético poderia fazer sobre este
trabalho, com foco em:

pontos fracos,

pressupostos discutíveis,

alternativas metodológicas.

Depois, para cada pergunta, sugere uma resposta
"razoável" que o candidato poderia dar.
\end{lstlisting}

Usa estas perguntas para treinar respostas em voz alta.
A vantagem é que a IA não se cansa de ser chata.

\textbf{Resultado esperado e validação:} Responde primeiro às perguntas sem ver a sugestão. Confronta cada resposta com a metodologia e os dados reais, mantendo limitações explícitas.

\textbf{Falhas possíveis:} Decorar uma defesa inventada ou ocultar uma limitação em vez de a explicar.
\end{promptlabbox}

\section{Uma ficha de evidência por afirmação importante}
Guarda título, autores, versão, URL/DOI, secção ou página, método, amostra ou modelos avaliados, resultado e limitações. Se só leste o abstract, declara-o: não podes auditar uma metodologia inteira a partir dele. Distingue o que os autores mediram da interpretação que estás a propor.

Uma referência existente pode continuar a ser inadequada: o artigo pode tratar outro domínio ou mostrar apenas um resultado num conjunto restrito. A pergunta não é só «esta fonte existe?», mas «sustenta esta frase, com este alcance?». O apêndice~\ref{ap:fontes} oferece um protocolo de leitura; a bibliografia reúne fontes concretas desta edição.
